%% ===========================================================================
%%  IEEE Transactions on Circuits and Systems for Video Technology (TCSVT)
%%  Special Issue: Towards Spatial Intelligence: 3D Understanding, Generation,
%%  and Reasoning from Real-World Videos  (submission deadline 2026-10-31)
%%
%%  Class:  IEEEtran.cls V1.8b (2015/08/26), vendored in vendor/IEEEtran/.
%%  Build:  make        (docker + TeX Live, see README.md)
%%
%%  RULE: every number in this file comes from a macro defined in numbers.tex,
%%        which is kept in step with docs/paper-numbers.md. Do not type a
%%        numeric result directly into the prose. The one exception is the
%%        per-scene table (Table III), transcribed from docs/muvod-comparison.md.
%% ===========================================================================
\documentclass[journal]{IEEEtran}

% --- IEEE-recommended packages ---------------------------------------------
\usepackage{cite}                 % IEEE reference style: [1], [2]-[5]
\usepackage[T1]{fontenc}
\usepackage[utf8]{inputenc}
\usepackage{graphicx}
\usepackage{amsmath,amssymb}
\usepackage{array}
\usepackage{booktabs}
\usepackage{multirow}
\usepackage[caption=false,font=footnotesize]{subfig}   % IEEEtran-recommended
\usepackage{xcolor}
\usepackage{tikz}
\usetikzlibrary{arrows.meta,positioning,fit,calc}
\usepackage{balance}              % even out the two columns on the last page
% hyperref must be configured before orcidlink, which loads it itself
\PassOptionsToPackage{hidelinks}{hyperref}
\usepackage{orcidlink}
\usepackage{hyperref}             % must stay last
\urlstyle{same}

\newcommand{\Jscore}{\ensuremath{\mathcal{J}}}
\newcommand{\Fscore}{\ensuremath{\mathcal{F}}}
\newcommand{\JF}{\ensuremath{\mathcal{J}\&\mathcal{F}^{3}}}
\newcommand{\JFpair}{\ensuremath{\mathcal{J}\&\mathcal{F}}}
\newcommand{\cini}{\ensuremath{c_{\mathrm{ini}}}}
\newcommand{\px}{\,px}

% --- drafting aids: set \draftfalse before submitting ----------------------
\newif\ifdraft\drafttrue
\ifdraft
  \newcommand{\todo}[1]{\textcolor{red}{\small$\blacktriangleright$\,\textsf{TODO: #1}}}
  \newcommand{\note}[1]{\textcolor{blue!70!black}{\small$\dagger$\,\textsf{#1}}}
\else
  \newcommand{\todo}[1]{}
  \newcommand{\note}[1]{}
\fi

%% ---------------------------------------------------------------------------
%% numbers.tex — 논문에 등장하는 모든 수치의 단일 출처.
%%
%% 이 파일은 docs/paper-numbers.md 와 1:1로 대응한다. 본문에서는 절대 숫자를
%% 직접 타이핑하지 말고 여기 매크로만 쓴다. 수치가 바뀌면 이 파일과
%% docs/paper-numbers.md 를 함께 고친다.
%% ---------------------------------------------------------------------------

%% §1 헤드라인 -----------------------------------------------------------------
\newcommand{\OursBasic}{91.015}
\newcommand{\OursComplete}{86.399}
\newcommand{\XMemBasic}{79.35}
\newcommand{\XMemComplete}{75.63}
\newcommand{\DeltaBasic}{+11.662}
\newcommand{\DeltaComplete}{+10.770}
\newcommand{\tBasic}{7.24}
\newcommand{\tComplete}{7.43}
\newcommand{\ciBasic}{[+8.71,\,+14.91]}
\newcommand{\ciComplete}{[+8.10,\,+13.68]}
\newcommand{\winLoss}{16/1}
\newcommand{\medianBasic}{+10.32}
\newcommand{\worstScene}{Welder}
\newcommand{\worstDelta}{\ensuremath{-0.73}}
\newcommand{\bestScene}{AlexaMeadeFacePaint}
\newcommand{\bestDelta}{+27.66}
\newcommand{\nScenes}{17}
\newcommand{\nBoot}{10\,000}

%% §2 분해 사다리 ---------------------------------------------------------------
\newcommand{\ladderXMem}{79.400}
\newcommand{\ladderTwoStage}{89.519}
\newcommand{\ladderXview}{90.034}
\newcommand{\ladderFsSeed}{90.064}
\newcommand{\ladderCrop}{91.015}
\newcommand{\gainFoundation}{+10.119}
\newcommand{\gainXview}{+0.515}
\newcommand{\gainFsSeed}{+0.030}
\newcommand{\gainCrop}{+0.951}
\newcommand{\gainTotal}{+11.615}
\newcommand{\gainOurs}{+1.50}
\newcommand{\pctFoundation}{87}

%% §3 선택 낙관 (LOSO) ----------------------------------------------------------
\newcommand{\losoFitted}{+1.053}
\newcommand{\losoHeldOut}{+1.008}
\newcommand{\losoOmega}{+0.045}
\newcommand{\losoOmegaPct}{4.3}
\newcommand{\losoAgree}{15}

%% §4 실행 시간·메모리 ----------------------------------------------------------
\newcommand{\runtimeTotal}{113}
\newcommand{\runtimeMedian}{302}
\newcommand{\runtimeMin}{100}
\newcommand{\runtimeMax}{1\,090}
\newcommand{\gpuMedian}{11.5}
\newcommand{\gpuMax}{20.8}
\newcommand{\gpuTotal}{48}

%% §5 오차 예산 ----------------------------------------------------------------
\newcommand{\totalGap}{8.985}
\newcommand{\oracleSeed}{+1.11}
\newcommand{\oracleRef}{+1.59}
\newcommand{\oracleOrdinary}{+7.87}
\newcommand{\nOrdinaryPairs}{1\,051}
\newcommand{\regionJ}{87.3}
\newcommand{\boundaryF}{94.8}

%% §6 격자 천장 ----------------------------------------------------------------
\newcommand{\gridNow}{288}
\newcommand{\ceilNow}{98.201}
\newcommand{\lossNow}{1.799}
\newcommand{\ceilMid}{99.340}
\newcommand{\lossMid}{0.660}
\newcommand{\ceilHigh}{99.784}
\newcommand{\lossHigh}{0.216}
\newcommand{\ceilSmall}{93.201}
\newcommand{\ceilMedium}{97.663}
\newcommand{\ceilLarge}{99.004}
\newcommand{\nRows}{1\,066}
\newcommand{\nFrames}{22\,220}

%% §7 1단계 음성 결과 -----------------------------------------------------------
\newcommand{\nQualPairs}{19}
\newcommand{\nQualPinhole}{16}
\newcommand{\gateDet}{9}
\newcommand{\gateCorr}{5}
\newcommand{\gateGeoA}{4}
\newcommand{\gateGeoB}{1}
\newcommand{\gateGeoC}{0}
\newcommand{\nGeoPairs}{16}
\newcommand{\unionOracle}{11}
\newcommand{\unionHalluc}{132/132}
\newcommand{\voteOracle}{+2.967}
\newcommand{\voteRule}{\ensuremath{-0.061}}
\newcommand{\nFzero}{41}
\newcommand{\nNzero}{623}
\newcommand{\asymmetry}{1{:}15}

%% §8 방법론 ------------------------------------------------------------------
\newcommand{\nPrereg}{22}
\newcommand{\nTests}{506}

%% 파생 표현 ------------------------------------------------------------------
\newcommand{\lossNowPctText}{20.0\% of the residual error}


\begin{document}

\title{Propagating a Single-View Mask Across a Static Camera Rig:\\
Attribution and Error Budget of a Foundation-Model Pipeline\\ for Multi-View Video Object Segmentation}
%% NOTE: working title. \note/\todo must never appear between \title and
%% \maketitle: IEEEtran's \maketitle issues \twocolumn[...] and any typeset
%% material before it pushes the whole title block onto the next page.

\author{Seongjin~Park\,\orcidlink{0000-0000-0000-0000}%
\thanks{Manuscript received XXX; revised XXX.}%
\thanks{S.~Park is with the Electronics and Telecommunications Research Institute
(ETRI), Daejeon, Republic of Korea (e-mail: seongjin.park@etri.re.kr).}%
\thanks{\todo{funding / project acknowledgement}}}

\markboth{IEEE Transactions on Circuits and Systems for Video Technology,~Vol.~XX, No.~X, XXX~2027}%
{Park: Propagating a Single-View Mask Across a Static Camera Rig}

\maketitle

% ===========================================================================
\begin{abstract}
A static multi-camera rig observes one scene from many viewpoints at once.
Given the mask of one object in one camera at one instant, the task is to
produce that object's mask in every other camera and every subsequent frame,
with no camera calibration, no category label, and no per-scene optimisation.
We study this task on MUVOD, the only public rig benchmark, and make three
contributions. First, a two-stage pipeline built on a promptable segmentation
foundation model reaches \OursBasic\ \JF\ (basic) and \OursComplete\
(complete), against \XMemBasic/\XMemComplete\ for the published baseline: a
paired gain of \DeltaBasic\ points ($t=\tBasic$, 95\% CI $\ciBasic$), winning
\winLoss\ scenes. Second, we decompose that margin rather than report it whole:
\gainFoundation\ points come from the foundation model inside an unchanged
two-stage structure, and \gainOurs\ from our own design---cross-view memory
during tracking, raw frame-0 seeding, and zoom-gated crop re-tracking---whose
selection optimism we bound at \losoOmega\ points by leave-one-scene-out.
Third, we account for the \totalGap\ points that remain: \oracleSeed\ in
seeding failures, \oracleRef\ in reference-camera re-prediction, and
\oracleOrdinary\ in ordinary mask quality on pairs that never fail. A
model-free measurement shows the tracker's \gridNow-cell mask grid costs only
\lossNow\ of those points. Five pre-registered attempts to repair degenerate
seeds---exemplar detection, dense correspondence, and three geometry-lifting
families---all fail, and we identify why: \nFzero\ pairs need repair against
\nNzero\ that must be left alone, a \asymmetry\ asymmetry that turns any
imperfect selector into a net loss.
\end{abstract}

\begin{IEEEkeywords}
Multi-view video object segmentation, camera rig, mask propagation, foundation
models, promptable segmentation, benchmark analysis, error budget,
pre-registration.
\end{IEEEkeywords}

% ===========================================================================
\section{Introduction}
\label{sec:intro}
\IEEEPARstart{A}{static} multi-camera rig---a light-field array, a volumetric
capture studio, an instrumented sports arena---observes one scene from many
viewpoints at the same instant. When a person marks one object in one of
those cameras, the natural expectation is that the mark should suffice: the
same object should be found in every other camera, and followed through time
in each. This paper is about that expectation. The input is one binary mask of
one object in one camera at one instant; the output is the mask of that object
in every camera and every subsequent frame. No camera calibration is provided,
no category label is attached to the object, and no per-scene training or
reconstruction is permitted.

The task sits between two well-studied problems and belongs to neither. It is
not video object segmentation: the cameras are not a temporal sequence, their
baselines are wide, and their index order carries no guarantee of visual
adjacency. It is not 3D segmentation in the sense of the NeRF and 3D Gaussian
Splatting (3DGS) literature: those methods assume calibrated poses and a
pre-trained scene representation, and the best of them are now saturated near
94\,mIoU on the NVOS benchmark \cite{sagonline,flashsplat,ludvig}. Treating the
views as a video and propagating with a memory-based tracker is the obvious
alternative, and its failure has been quantified: the segment-correspondence
AUPRC of SAM~2 video propagation on ScanNet++ falls from 61.9 at
$0$--$45^{\circ}$ view difference to 17.0 at $135$--$180^{\circ}$
\cite{segmast3r}. The 2026 generation of uncalibrated promptable methods
\cite{mvsam,am3,g2tam} is trained on indoor, object-centric data and has not
been evaluated on rigs, on environment objects such as wall and floor
fragments, or on fisheye cameras.

Only one public benchmark poses the rig problem directly. MUVOD \cite{muvod}
provides \nScenes\ scenes captured by \nCamMin--\nCamMax\ synchronised
cameras, with \nInstances\ object instances over \nCategories\ categories,
\nFisheye\ hemispherical fisheye rigs among them, and a published baseline that
applies a video tracker along both the spatial and the temporal axis. At the
time of writing no other method has been evaluated on it. This has a
consequence for what an honest paper on the task can claim: the first strong
number on MUVOD will be informative mainly in proportion to how carefully it is
taken apart.

We therefore build the simplest pipeline that could plausibly work---the
benchmark's own two-stage structure with the video tracker replaced by a
promptable segmentation foundation model, SAM~3 \cite{sam3}---add three small
design elements, and then spend most of the paper measuring rather than
claiming. The pipeline reaches \OursBasic\ \JF\ on the \emph{basic} protocol
and \OursComplete\ on \emph{complete}, against \XMemBasic\ and \XMemComplete\
for the published XMem baseline (Table~\ref{tab:headline}). Reported as a
single margin, that is \DeltaBasic\ points; reported that way it would also be
misleading, because \pctFoundation\% of it is attributable to the foundation
model inside a structure we did not change, and only \gainOurs\ points to our
design (Table~\ref{tab:ladder}). We say so, and we bound the optimism of the
one design choice made by looking at the scored scenes at \losoOmega\ points
(Table~\ref{tab:loso}).

The second half of the paper asks where the remaining \totalGap\ points are.
An oracle decomposition (Table~\ref{tab:budget}) attributes \oracleSeed\ of
them to the pairs whose cross-view seed fails, \oracleRef\ to the ceiling of the
reference camera itself, and \oracleOrdinary\ to ordinary, unspectacular mask
quality on the \nOrdinaryPairs\ pairs that never fail. A model-free
measurement---ground truth passed through the tracker's own low-resolution mask
grid---shows that the grid accounts for only \lossNow\ points
(Table~\ref{tab:grid}), so the large bucket is genuine model error rather than
a resolution artefact. Finally we report five pre-registered attempts to
repair the failing seeds by exemplar detection, dense correspondence, and three
geometry-lifting routes (Table~\ref{tab:negative}). All five fail their
registered gates, a ground-truth-free vote over them loses \voteRule\ points
where an oracle selector would gain \voteOracle, and the reason is structural:
there are \nFzero\ pairs to repair and \nNzero\ to protect.

\noindent\textbf{Contributions.}
\begin{itemize}
  \item A two-stage rig pipeline that sets the first foundation-model result on
        MUVOD (\OursBasic/\OursComplete\ \JF), with a paired scene-cluster
        analysis against the published baseline and a \emph{decomposition} of
        the margin that separates the foundation model (\gainFoundation) from
        our design (\gainOurs), plus a leave-one-scene-out bound on selection
        optimism ($\Omega=\losoOmega$).
  \item An error budget for the remaining \totalGap\ points, including a
        model-free measurement of the mask-grid ceiling (\lossNow\ points) and
        the observation that the non-reference cameras' deficit is constant
        from the first frame rather than accumulated drift.
  \item Five pre-registered negative results on seed repair, the union-oracle
        versus hallucination trade-off that underlies them, and the
        \asymmetry\ repair-to-protect asymmetry as the quantity future work on
        this task must report against.
  \item A reconstruction of the MUVOD reference-camera protocol, which the
        benchmark does not publish, from three independent lines of evidence.
\end{itemize}

% ===========================================================================
\section{Related Work}
\label{sec:related}

\subsection{Segmentation on a reconstructed 3D representation}
SPIn-NeRF \cite{spinnerf} and SA3D \cite{sa3d} formalised ``one-view prompt
$\rightarrow$ all-view masks'' on top of a per-scene radiance field, and the
representation moved to 3DGS with Gaussian Grouping \cite{gaussiangrouping} and
its successors. The current generation lifts masks onto a \emph{pre-trained}
3DGS without further optimisation---FlashSplat \cite{flashsplat}, LUDVIG
\cite{ludvig}, SAGOnline \cite{sagonline}---and has pushed NVOS mIoU to the
point where its own authors describe the remaining differences between methods
as converging on the 2D foundation model's ceiling \cite{sagonline}. Every
member of this family assumes COLMAP poses \cite{colmap} and a reconstruction
of sufficient quality; LUDVIG and SANeRF-HQ state explicitly that the
reconstruction bounds the segmentation. Sec.~\ref{sec:negative} tests exactly
this bridge on rigs of \nCamMin--\nCamMax\ cameras and finds that on
\gsCams\ of \gsCams\ cameras the rendered depth is biased towards the camera,
with the bias \emph{anti}-correlated with training-view PSNR.

\subsection{Views as a video}
Ordering the cameras and propagating a first-frame mask with a video object
segmentation tracker is the structure of the MUVOD baseline \cite{muvod}, of
OSTRA \cite{ostra}, and of the pre-processing stage of Gaussian Grouping; the
memory-based trackers XMem \cite{xmem}, DEVA \cite{deva}, Cutie \cite{cutie},
and SAM~2 \cite{sam2} are the usual engines. WildSeg3D \cite{wildseg3d} adds
MASt3R alignment \cite{mast3r} to choose the order. The approach inherits the
tracker's prior that consecutive frames are nearly aligned, and it degrades
accordingly at wide baselines \cite{segmast3r}. Our two-stage baseline belongs
to this family, and we say so plainly: its cross-view stage is a single
bidirectional pass of a SAM~3 tracker over the cameras' first frames in index
order. What we add on top (Sec.~\ref{sec:method}) is small, and we measure it
as small.

\subsection{Correspondence, uncalibrated promptable methods}
Training-free feature correspondence---Matcher \cite{matcher}, PerSAM
\cite{persam}---turns a reference mask into prompts for a query image through
DINOv2 \cite{dinov2} or SAM \cite{sam} features; its sensitivity to viewpoint change is
documented \cite{probing3d}. Feed-forward geometry models
\cite{dust3r,mast3r,vggt,pi3} now provide poses and point maps without
calibration, and a 2026 family injects them into promptable segmenters:
MV-SAM \cite{mvsam}, 3AM \cite{am3}, and G$^{2}$TAM \cite{g2tam}, alongside
learned segment matching such as SegMASt3R \cite{segmast3r}. These are the
closest existing designs to our input format. None reports MUVOD, none
evaluates environment objects, and none has been run on fisheye rigs; that is
why they do not appear as columns in Table~\ref{tab:headline}, a limitation we
return to in Sec.~\ref{sec:limits}. We do, however, evaluate the training-free
ingredients they share---DINOv2 correspondence, $\pi^{3}$ point maps, and SAM~3
exemplar detection---as seed generators in Sec.~\ref{sec:negative}.

\subsection{Rig benchmarks}
MUVOD \cite{muvod} is, to our knowledge, the only public benchmark with
synchronised multi-camera rigs, per-instance masks, and a fixed protocol. Its
3D benchmark excludes the \nFisheye\ fisheye scenes and one further scene
whose reconstruction was judged too poor, and it uses camera poses; its video
object segmentation benchmark, which we use, includes all \nScenes\ scenes and
provides no calibration. Its reference camera per scene is not published,
which we address in Sec.~\ref{sec:cini}.

% ===========================================================================
\section{Problem Setting and Protocol}
\label{sec:protocol}

\subsection{Task}
A rig of $N$ synchronised cameras $c_{1},\dots,c_{N}$ records $T+1$ frames
each. One camera, the reference \cini, carries a set of binary object masks
$\{M_{o}\}$ at frame $t=0$. For every object $o$, every camera $c$, and every
frame $t\in\{0,\dots,T\}$ the method must output a mask $\hat{M}_{o,c,t}$.
Intrinsics and extrinsics are not provided and are not estimated for use in
the pipeline; no scene-specific training or optimisation is performed; the
segmenter is a frozen pre-trained model. All frames of all cameras may be
read---the benchmark's own annotation procedure and its baseline both assume
this.

\subsection{The MUVOD metric}
MUVOD scores \nAnnotCams\ annotated cameras per scene: \cini\ and two others.
For each scored camera, region similarity \Jscore\ (intersection over union)
and boundary accuracy \Fscore\ are computed per object over the
\nFramesPerCam\ annotated frames with the DAVIS definitions \cite{davis2016};
the per-camera score is the mean over objects, the scene score is the mean over
the three cameras, and the benchmark score is the mean over \nScenes\ scenes.
We write the result \JF. Two object sets are defined: \emph{basic} counts only
objects visible in the reference frame, \emph{complete} counts every labelled
object. The two coincide for a scene exactly when the reference frame already
contains all of that scene's labelled objects, a property we use below. Our
scoring code is the DAVIS implementation, unchanged; two choices the benchmark
leaves open we fix conservatively---an object absent from a camera's ground
truth is not scored for that camera (the alternative awards a perfect score to
an empty prediction), and all \nFramesPerCam\ frames are scored rather than
the DAVIS convention of dropping the first and last.

\subsection{Recovering the reference camera}
\label{sec:cini}
The benchmark describes \cini\ as a camera ``near the centre of the rig'' but
does not name it, and the \emph{basic} object set depends on it. We determined
it per scene from three independent sources. (A) The distributed index records
the rig geometry; on linear, arc, and planar rigs the middle-numbered annotated
camera is the geometric centre. (B) For the six scenes whose camera numbering
is not spatial---four hemispherical rigs and two two-row planar rigs---we read
the original datasets' calibration files and chose the annotated camera closest
to the rig centre in metric coordinates. (C) The published tables themselves
constrain the answer: in eight scenes the basic and complete scores are equal,
which means the reference frame contains every labelled object, and in three of
those scenes exactly one annotated camera satisfies that condition in the
ground truth. In those three scenes (C) agrees with (A) or (B). One scene, Dog,
remains ambiguous: (C) excludes one of its three annotated cameras, the
remaining two lie on the same ring of the hemisphere, and we adopt the one
marginally nearer the axis. A benchmark run with the other candidate is
reported alongside. The choice matters: the reference rule of our own earlier
runs---the camera with the most objects at frame 0---selects a different camera
from \cini\ in \nCiniMismatch\ of \nScenes\ scenes, typically one at the edge
of the rig.

% ===========================================================================
\section{Method}
\label{sec:method}

\begin{figure}[!t]
\centering
\begin{tikzpicture}[
  font=\footnotesize,
  box/.style={draw, rounded corners=2pt, minimum height=1.6em, inner sep=3pt, align=center},
  stage/.style={draw, dashed, rounded corners=3pt, inner sep=5pt},
  arr/.style={-{Latex[length=1.6mm]}, thick},
  node distance=4mm and 6mm]
  % stage 1
  \node[box] (f0) {first frames\\$c_{1}\dots c_{N}$};
  \node[box, right=of f0] (pv) {pseudo-video\\index order};
  \node[box, right=of pv] (s1) {SAM~3 tracker\\fwd + bwd};
  \node[box, right=of s1] (seed) {seeds\\$\hat M_{o,c,0}$};
  \node[above=5mm of f0, anchor=south west, xshift=-3mm] {\textbf{Stage 1} (cross-view, one instant)};
  \draw[arr] (f0) -- (pv);
  \draw[arr] (pv) -- (s1);
  \draw[arr] (s1) -- (seed);
  \node[box, above=5mm of s1, fill=black!6] (gt) {$M_{o}$ in \cini};
  \draw[arr] (gt) -- (s1);
  % stage 2
  \node[box, below=9mm of f0] (t1) {camera $c$\\$t=1\dots T$};
  \node[box, right=of t1] (trk) {SAM~3 tracker\\per camera};
  \node[box, right=of trk] (out) {$\hat M_{o,c,t}$};
  \node[box, right=of out, fill=black!6] (crop) {crop re-track\\seed $<\cropAreaMax$\px};
  \node[below=1mm of trk, xshift=6mm] {\textbf{Stage 2} (temporal, per camera)};
  \draw[arr] (t1) -- (trk);
  \draw[arr] (trk) -- (out);
  \draw[arr] (out) -- (crop);
  \draw[arr] (seed) |- ($(trk.north)+(0,3mm)$) -- (trk.north);
  % cross-view memory
  \node[box, below=of trk, fill=black!6, yshift=-4mm] (nb) {neighbour $c\pm1$ memory\\gate, pointer, row \xviewRow};
  \draw[arr] (nb) -- (trk);
\end{tikzpicture}
\caption{The pipeline. Stage~1 propagates the reference masks across the
cameras' first frames once; Stage~2 tracks each camera forward in time from
its seed. Shaded blocks are the elements we add to the baseline structure:
cross-view memory during tracking (Sec.~\ref{sec:xview}) and crop re-tracking
of small objects (Sec.~\ref{sec:crop}); the seed is also written unchanged as
the frame-0 output (Sec.~\ref{sec:fs}).}
\label{fig:pipeline}
\end{figure}

\subsection{A two-stage baseline}
\label{sec:twostage}
Fig.~\ref{fig:pipeline} shows the structure. \emph{Stage~1} arranges the first
frames of all $N$ cameras in camera-index order as a pseudo-video, gives the
reference masks $\{M_{o}\}$ as the only prompt on the frame of \cini, and runs
the SAM~3 video tracker once forward to the end of the sequence and once
backward to its start. The result is one seed mask $\hat{M}_{o,c,0}$ per
(camera, object). \emph{Stage~2} opens one temporal session per camera, prompts
it with that camera's seeds at $t=0$, and tracks forward to $t=T$. The segmenter
is SAM~3 \cite{sam3}, whose tracker inherits the memory-attention design of
SAM~2 \cite{sam2}: a bank of memory tokens from recent frames, each carrying a
learned temporal-position encoding, plus per-object pointer vectors. It runs at
a \modelInput-pixel input and predicts masks on a \gridNow-cell low-resolution
grid before upsampling, a fact that matters in Sec.~\ref{sec:grid}.

This is structurally the MUVOD baseline. The benchmark authors describe theirs
as a ``simple two-stage baseline'' that applies XMem along the spatial axis and
then along the temporal axis, and note its limitation as being ``inherently
based on mask propagation'' without multi-view understanding. Ours is that
structure with the tracker exchanged. We verified the claim that our control
configuration contains nothing else by byte comparison: with cross-view memory
disabled, Stage~2 writes \nBytePNG\ PNG masks identical to those of an
independent single-view package that has no cross-view code at all.

\subsection{Cross-view memory during tracking}
\label{sec:xview}
In the baseline structure the cameras never communicate after Stage~1. We let
them communicate through the tracker's own memory: at frame $t$, camera $c$'s
memory bank additionally receives the memory token and object pointer of its
index neighbours within a window $W=\xviewWindow$---the previous camera at
frame $t$ and the next camera at frame $t-1$, so that each read refers to a
frame that has already been computed. Three conditions govern the neighbour
tokens. They pass the same quality gate as the camera's own memory (a
low-confidence neighbour frame is dropped exactly as a low-confidence own frame
would be); the neighbour's object pointer is appended after the camera's own
pointers; and the neighbour token is assigned the temporal-position encoding of
row \xviewRow, i.e.\ ``\xviewRow\ frames ago'', rather than row 0.

The last choice is grounded in a measurement. Comparing pixel change between
adjacent cameras at the same instant with pixel change between the same camera
at $t$ and $t-k$, one camera step exceeds $k=\viewStepFrames$ frames of change
in \viewStepPct\% of camera pairs on every scene; the tracker's temporal
encoding has only seven rows, so the neighbour is placed as far back as the
encoding allows short of the conditioning row. The curve over the shift is not
monotonic: it peaks at row \xviewRow\ and falls again at row 6. Window size
saturates immediately---on our development split $W=2$, $4$ and $6$ add
nothing over $W=1$---and the gate, pointer, and shift, each individually
positive, overlap rather than add when combined.

We report the history of this component honestly because it bears on how much
to trust it. Its first pre-registered test, at our original edge-of-rig
reference cameras, found no effect (paired $\Delta$ on the non-reference
cameras that receive a neighbour: \firstTestDelta, 95\% CI \firstTestCI). A
direction ablation likewise found no mode passing its registered criterion.
Only after the protocol moved to the benchmark's central \cini\ did a
re-registered test pass: the \nNeighbourless\ cameras that previously had no
neighbour gained \ciniTestDelta\ (CI \ciniTestCI), and the benchmark moved by
\ciniTestScenes\ points with no scene falling below the registered floor. The
mechanism did not change; the number of cameras it could act on did.

\subsection{Raw frame-0 seeding}
\label{sec:fs}
Stage~2 re-predicts frame~0 from the seed it is prompted with, and the
re-prediction is slightly worse than the seed. We write the seed itself as the
frame-0 output. The effect is \gainFsSeed\ points; it is reported because it is
one rung of the ladder in Table~\ref{tab:ladder}, not because it matters.

\subsection{Zoom-gated crop re-tracking}
\label{sec:crop}
Small objects lose most on the \gridNow-cell grid. For every (camera, object)
whose seed covers fewer than $\cropAreaMax$ pixels we open a second, single-object
tracking session inside a square window centred on the seed's bounding box,
with side $s=\mathrm{clamp}(\cropScale\cdot\max(w,h),\ \cropMinSide,\
\min(H,W))$, resized to the model's \modelInput-pixel input; the effective zoom
is $\modelInput/s$. The session reads no neighbour memory and its result is
pasted back into the full-frame output for frames $1$ to $T$. The seed rule
uses the seed, not the ground truth, so it is available at inference. On the
reference camera, whose seeds are ground truth, the \nCropRefSmall\ objects
under $2\times10^{3}$\px\ rise from a mean frame-1--20 \Jscore\ of
\cropRefSmallBefore\ to \cropRefSmallAfter, which is the mechanism working as
intended.

A window can also contain a second object of the same appearance, and the
tracker then follows the wrong one from the first frame. A ground-truth-free
gate reverts a pair to the full-frame result when any of three conditions
holds: the crop-tracked mask at frame~1 overlaps the seed with IoU below
\gateIoU; the zoom is below \gateZoom, so the crop gains no resolution; or the
object touches the window edge in \gateBorder\ or more of the 20 tracked frames.
Across the benchmark \nCropTargets\ pairs qualify and \nCropReverted\ are
reverted. The thresholds were chosen on the scored scenes; Sec.~\ref{sec:loso}
measures what that cost.

\subsection{Pre-registration and controls}
\label{sec:prereg}
Every experiment in this paper was pre-registered: its hypothesis, population,
acceptance rule, and predicted outcome were written and committed to the
repository before the run, and no rule was changed after the data were seen.
Where a rule was later re-registered we say so (Sec.~\ref{sec:xview}). There
are \nPrereg\ such documents. The default code path is guarded by \nTests\ CPU
tests and by byte-level identity checks: cameras that receive no neighbour
change by exactly zero; the \nByteNontarget\ masks of pairs not touched by crop
re-tracking are byte-identical to the previous configuration; the
\nByteSeeds\ reference-camera seeds are byte-identical across every Stage-1
variant in Sec.~\ref{sec:negative}.

% ===========================================================================
\section{Experiments}
\label{sec:exp}

\subsection{Setup}
All runs use one NVIDIA RTX~6000 Ada (\gpuTotal~GiB) in a Docker container,
fp32, deterministic. The benchmark is the full MUVOD video set: \nScenes\
scenes, \nAnnotCams\ scored cameras per scene, \nFramesPerCam\ frames per
camera, \nPairsBasic\ (camera, object) pairs in the basic set and
\nPairsComplete\ in the complete set; the \nPairsUnreachable\ pairs in the
difference are objects absent from the reference frame, which no propagation
method can reach. Paired statistics use a scene-cluster bootstrap with
\nBoot\ resamples and a fixed seed. The XMem baseline numbers are transcribed
from the benchmark paper; we did not re-run it (Sec.~\ref{sec:limits}).

\begin{table}[!t]
\caption{Comparison with the published MUVOD baseline. \JF\ (\%) over
\nScenes\ scenes. $\Delta$ and CI are paired per scene (cluster bootstrap,
\nBoot\ resamples).}
\label{tab:headline}
\centering
\begin{tabular}{lcccc}
\toprule
 & Ours & XMem~\cite{muvod} & $\Delta$ & 95\% CI \\
\midrule
basic    & \textbf{\OursBasic}    & \XMemBasic    & \textbf{\DeltaBasic}    & $\ciBasic$ \\
complete & \textbf{\OursComplete} & \XMemComplete & \textbf{\DeltaComplete} & $\ciComplete$ \\
\bottomrule
\end{tabular}
\end{table}

\subsection{Comparison with the published baseline}
Table~\ref{tab:headline} gives the benchmark result. On the basic protocol the
paired difference is \DeltaBasic\ points ($t=\tBasic$), on complete
\DeltaComplete\ ($t=\tComplete$); both intervals exclude zero by a wide margin.
The pipeline wins \winLoss\ scenes, with a median gain of \medianBasic; the
largest gain is on \bestScene\ (\bestDelta) and the one loss is \worstScene\
(\worstDelta). Table~\ref{tab:scenes} lists every scene. Scenes whose basic and
complete scores coincide are those whose reference frame contains every
labelled object, the same property that fixed \cini\ in Sec.~\ref{sec:cini}.

\begin{table*}[!t]
\caption{Per-scene results, \JF\ (\%). XMem values are transcribed from the
MUVOD paper. Scenes are ordered by basic $\Delta$.}
\label{tab:scenes}
\centering
\begin{tabular}{l ccc ccc}
\toprule
 & \multicolumn{3}{c}{basic} & \multicolumn{3}{c}{complete} \\
\cmidrule(lr){2-4}\cmidrule(lr){5-7}
Scene & Ours & XMem & $\Delta$ & Ours & XMem & $\Delta$ \\
\midrule
AlexaMeadeFacePaint & 92.561 & 64.9 & $+27.66$ & 83.987 & 60.4 & $+23.59$ \\
CBABasketball       & 88.564 & 66.4 & $+22.16$ & 87.021 & 65.6 & $+21.42$ \\
Breakfast           & 87.822 & 68.8 & $+19.02$ & 82.774 & 64.1 & $+18.67$ \\
MATF                & 85.397 & 69.1 & $+16.30$ & 73.083 & 59.3 & $+13.78$ \\
Dog                 & 88.348 & 75.5 & $+12.85$ & 77.896 & 67.2 & $+10.70$ \\
Barn                & 89.744 & 77.3 & $+12.44$ & 89.744 & 77.3 & $+12.44$ \\
Blocks              & 91.285 & 80.5 & $+10.78$ & 82.210 & 73.2 & $+9.01$ \\
Painter             & 92.645 & 82.2 & $+10.45$ & 85.957 & 76.1 & $+9.86$ \\
CoffeeMartini       & 94.621 & 84.3 & $+10.32$ & 85.620 & 76.7 & $+8.92$ \\
PoznanStreet        & 90.363 & 80.2 & $+10.16$ & 90.363 & 80.2 & $+10.16$ \\
MartialArts         & 93.785 & 83.7 & $+10.09$ & 93.785 & 83.7 & $+10.09$ \\
Carpark             & 94.495 & 85.3 & $+9.19$  & 94.495 & 85.3 & $+9.19$ \\
Fencing             & 93.197 & 85.7 & $+7.50$  & 93.197 & 85.7 & $+7.50$ \\
FlameSteak          & 92.129 & 84.8 & $+7.33$  & 76.359 & 70.6 & $+5.76$ \\
AlexaMeadeExhibit   & 90.047 & 82.8 & $+7.25$  & 90.047 & 82.8 & $+7.25$ \\
Frog                & 97.778 & 92.3 & $+5.48$  & 97.778 & 92.3 & $+5.48$ \\
Welder              & 84.467 & 85.2 & $-0.73$  & 84.467 & 85.2 & $-0.73$ \\
\midrule
Mean (\nScenes)     & \textbf{\OursBasic} & \XMemBasic & $\mathbf{\DeltaBasic}$ & \textbf{\OursComplete} & \XMemComplete & $\mathbf{\DeltaComplete}$ \\
\bottomrule
\end{tabular}
\end{table*}

\begin{table}[!t]
\caption{Where the margin comes from. \JF\ (\%), basic protocol. Each row
adds one element to the row above; the same scoring is applied throughout.}
\label{tab:ladder}
\centering
\begin{tabular}{lcc}
\toprule
Configuration & \JF & Gain \\
\midrule
XMem baseline~\cite{muvod}                       & \ladderXMem     & --- \\
Two-stage structure, SAM~3 tracker (Sec.~\ref{sec:twostage}) & \ladderTwoStage & \textbf{\gainFoundation} \\
$+$ cross-view memory (Sec.~\ref{sec:xview})     & \ladderXview    & \gainXview \\
$+$ raw frame-0 seed (Sec.~\ref{sec:fs})         & \ladderFsSeed   & \gainFsSeed \\
$+$ zoom-gated crop re-tracking (Sec.~\ref{sec:crop}) & \ladderCrop & \textbf{\gainCrop} \\
\midrule
Total                                            &                 & \textbf{\gainTotal} \\
\bottomrule
\end{tabular}
\end{table}

\subsection{Decomposition of the margin}
\label{sec:ladder}
Table~\ref{tab:ladder} climbs from the published baseline to our result one
element at a time. The first step---replacing the tracker inside an unchanged
two-stage structure---is worth \gainFoundation\ points, \pctFoundation\% of
the total. The three elements we designed contribute \gainXview, \gainFsSeed,
and \gainCrop, together \gainOurs. This is the paper's central number, and
two things must be said about it. First, the \gainFoundation\ step cannot be
split further into ``foundation model'' and ``two-stage structure'': we never
ran XMem inside our pipeline, so the structural equivalence with the benchmark
baseline is argued from the benchmark's description, not measured. Second, and
for the same reason, the converse claim---that any modern tracker would produce
this number---is equally unmeasured; no foundation-model method has been
evaluated on MUVOD by anyone. What is established is the size of our own
design's share, and that it is small.

\begin{table}[!t]
\caption{Leave-one-scene-out estimate of the selection optimism in the crop
gate. Gain in \JF\ (\%) over the configuration on which the choice was
actually made.}
\label{tab:loso}
\centering
\begin{tabular}{lc}
\toprule
 & Gain \\
\midrule
Rule chosen on all \nScenes\ scenes (fitted)   & \losoFitted \\
Rule refitted on 16, read on the held-out scene & \losoHeldOut \\
\midrule
Selection optimism $\Omega$                    & \textbf{\losoOmega} (\losoOmegaPct\%) \\
\bottomrule
\end{tabular}
\end{table}

\subsection{How much of the gain is selection optimism?}
\label{sec:loso}
The crop area threshold and the gate variant of Sec.~\ref{sec:crop} were
chosen among six candidate arms by looking at all \nScenes\ scored scenes.
We estimate the resulting optimism by leave-one-scene-out: for each scene the
same selection rule---largest mean gain subject to no scene falling more than
0.3 points---is re-applied to the other 16, and the chosen arm's gain is read
on the held-out scene. Table~\ref{tab:loso} gives the result. The fitted gain
is \losoFitted, the held-out gain \losoHeldOut, and their difference
$\Omega=\losoOmega$ points, \losoOmegaPct\% of the effect; \losoAgree\ of
\nScenes\ folds select the same arm. The objection ``tuned on the test set''
is valid; its magnitude is \losoOmega\ points.

\begin{table}[!t]
\caption{Cost. One RTX~6000 Ada (\gpuTotal~GiB), Docker, \nScenes\ scenes,
one configuration.}
\label{tab:runtime}
\centering
\begin{tabular}{lc}
\toprule
 & Value \\
\midrule
Total, \nScenes\ scenes    & \runtimeTotal~min \\
Per scene, median          & \runtimeMedian~s \\
Per scene, min--max        & \runtimeMin--\runtimeMax~s \\
Peak GPU memory, median    & \gpuMedian~GiB \\
Peak GPU memory, max       & \gpuMax~GiB \\
\bottomrule
\end{tabular}
\end{table}

\subsection{Runtime and memory}
Table~\ref{tab:runtime}. Time scales with the number of objects and the frame
resolution; there is no per-scene optimisation and no reconstruction, which
is the practical difference from the 3DGS family, where a reconstruction of
several minutes precedes any segmentation.

% ===========================================================================
\section{What Limits the Remaining Error}
\label{sec:budget}

\begin{table}[!t]
\caption{Oracle decomposition of the gap to a perfect score, basic protocol.
Each row replaces one component by ground truth and reports the gain.}
\label{tab:budget}
\centering
\begin{tabular}{lc}
\toprule
Component made perfect & Gain (\JF) \\
\midrule
Total gap ($100-\OursBasic$)                            & \totalGap \\
\midrule
Every frame-0 seeding failure                           & \oracleSeed \\
Reference camera \cini\ (all frames)                    & \oracleRef \\
The \nOrdinaryPairs\ pairs whose seed does not fail     & \textbf{\oracleOrdinary} \\
\bottomrule
\end{tabular}
\end{table}

\subsection{Oracle decomposition}
Table~\ref{tab:budget} asks, for each component, how much the benchmark score
would rise if that component alone were perfect. The pairs whose seed fails
at frame~0 (\JFpair${}<0.5$) are few, and repairing all of them is worth
\oracleSeed\ points. The reference camera---which starts from ground truth---has
a ceiling of \oracleRef. The large bucket, \oracleOrdinary\ points, is spread
over the \nOrdinaryPairs\ pairs whose seed is sound: \nGoodPairs\ of the
\nPairsBasic\ pairs already exceed 0.95, and the deficit is neither dominated by
a few catastrophes nor uniform. Region \Jscore\ (\regionJ) is worse than
boundary \Fscore\ (\boundaryF), though MUVOD's boundary tolerance of
\boundaryTol\px\ makes \Fscore\ lenient.

\subsection{A constant deficit, not drift}
The reference camera is prompted with ground truth, yet its per-frame \Jscore\
is \refJfZero\ at $t=0$, \refJfOne\ at $t=1$, and then flat
(\refJfFive, \refJfTen, \refJfTwenty\ at $t=5,10,20$). The drop occurs at the
first tracked frame and does not accumulate. The non-reference cameras show the
same shape at a lower level: their \JFpair\ is \nonrefJFzero\ at
$t=0$, \nonrefJFfive\ at $t=5$, \nonrefJFtwenty\ at $t=20$, against
\refJFzero, \refJFfive, \refJFtwenty\ on the reference camera. The deficit is
present from the first frame and constant thereafter. This rules out the
intuitive story that cross-view errors compound over time, and points at the
seed and at the tracker's per-frame quality rather than at temporal drift.

\begin{table}[!t]
\caption{Model-free ceiling of the tracker's mask grid: ground truth passed
through the same grid and scored against itself (\nRows\ pairs, \nFrames\
frames).}
\label{tab:grid}
\centering
\begin{tabular}{lccc}
\toprule
Grid & Ceiling \JF & Unreachable & Share of \totalGap \\
\midrule
\gridNow\ (ours) & \ceilNow  & \textbf{\lossNow} & 20.0\% \\
576              & \ceilMid  & \lossMid          & 7.4\% \\
1008             & \ceilHigh & \lossHigh         & 2.4\% \\
\bottomrule
\end{tabular}
\end{table}

\subsection{The mask grid is not the bottleneck}
\label{sec:grid}
The largest single gain in this paper came from a resolution device
(Sec.~\ref{sec:crop}), which invites the hypothesis that resolution is the
dominant residual. We tested it without a model: the ground-truth masks are
downsampled to the tracker's \gridNow-cell grid, upsampled back, and scored
against the originals under the benchmark protocol (Table~\ref{tab:grid}). The
grid's ceiling is \ceilNow, so it can account for at most \lossNow\ of the
\totalGap\ points; a 576-cell grid would leave \lossMid\ and a full-resolution
grid \lossHigh. The loss depends on object size alone---\ceilSmall\ for the
\nSmallRows\ pairs under $2\times10^{3}$\px, \ceilMedium\ for the
\nMediumRows\ between $2\times10^{3}$ and $10^{4}$, \ceilLarge\ for the
\nLargeRows\ above---and not on rig type or camera role. Crop re-tracking
already recovers \gainCrop\ of the \lossNow. The hypothesis is therefore
rejected: the \oracleOrdinary-point bucket of Table~\ref{tab:budget} is genuine
model error. Because the measurement uses no model, it is a property of the
benchmark and applies to any tracker with the same output grid.

% ===========================================================================
\section{Five Ways to Repair a Failing Seed, and Why None Works}
\label{sec:negative}

The oracle in Table~\ref{tab:budget} puts a bound of \oracleSeed\ points on
seed repair, but the failing seeds are also the most visible errors---an object
placed on the wrong thing, or missing altogether---and the literature of
Sec.~\ref{sec:related} offers several training-free ways to produce a better
one. We tested five, each pre-registered with its acceptance gate, and report
all five.

\subsection{Population}
Stage~1 is evaluated on its own output, the seed at $t=0$, on the two
non-reference scored cameras of every scene: \nRowsStage\ (camera, object) rows,
of which \nGtPresent\ have ground truth---\nFzero\ \emph{failing} (raw seed
\Jscore$<\voteTau$; we call this set F0) and \nNzero\ \emph{sound} (N0)---and
\nGtEmpty\ have no ground truth in that camera, where any non-empty candidate
is a hallucination. The existence gates are read on the \nQualPairs\ failing
pairs with ground truth of at least $2\times10^{3}$\px\ (\nQualPinhole\
pinhole, \nQualHemi\ hemisphere), and the geometry arms on the \nGeoPairs\
pinhole pairs among them, because none of the geometry models handles the raw
fisheye frames of the hemispherical rigs.

\begin{table}[!t]
\caption{Pre-registered seed generators. Gate: number of qualifying failing
pairs on which a candidate with IoU$>0.5$ against ground truth exists. The
vote's gate is a paired improvement whose CI excludes zero \emph{and} reaches
$+0.3$ benchmark points.}
\label{tab:negative}
\centering
\footnotesize
\begin{tabular}{@{}llcc@{}}
\toprule
Generator & Position source & Gate & Result \\
\midrule
G-det   & SAM~3 exemplars        & $\geq 11$/\nQualPairs & \gateDet/\nQualPairs \\
G-corr  & DINOv2 matching        & $\geq 11$/\nQualPairs & \gateCorr/\nQualPairs \\
G-geo-a & SfM $+$ dense MVS      & $\geq 8$/\nGeoPairs   & \gateGeoA/\nGeoPairs \\
G-geo-b & $\pi^{3}$ point maps   & $\geq 8$/\nGeoPairs   & \gateGeoB/\nGeoPairs \\
G-gs    & 3DGS depth             & $\geq 8$/\nGeoPairs   & \gateGeoC/\nGeoPairs \\
\midrule
SeedVote & all above $+$ order   & CI${}>0$, ${\geq}+0.3$ & \voteRule \\
\bottomrule
\end{tabular}
\end{table}

\subsection{The generators}
Table~\ref{tab:negative} lists them. In every arm the mask is produced by the
frozen SAM~3; the arms differ in how the object's \emph{position} in the target
camera is obtained. \emph{G-det} injects the reference mask as an exemplar
prompt into SAM~3's detector on each target frame and takes the union of the
box- and point-exemplar candidates. \emph{G-corr} builds a DINOv2-L prototype
from the patch tokens under the reference mask, computes its cosine map on the
target frame, and prompts a single-frame SAM~3 session at the top peaks.
\emph{G-geo-a} runs COLMAP structure-from-motion on the first frames, densifies
with multi-view stereo, lifts the reference mask to 3D, and reprojects it.
\emph{G-geo-b} replaces the reconstruction by $\pi^{3}$ feed-forward point maps
\cite{pi3}. \emph{G-gs} trains a 3DGS \cite{gaussiansplatting} on the first
frames for \gsIter\ iterations, renders $\alpha$-normalised depth, and uses the
same lifting.

None passes. The interesting part is \emph{how} each fails, because the
failures are not the same. G-det finds a correct candidate on \gateDet\ of
\nQualPairs\ failing pairs but also produces at least one candidate on
\detHalluc\ of the ground-truth-empty rows. G-corr works on sound pairs---a
correct candidate exists on \corrSoundExist\% of N0 with median IoU
\corrSoundIoU---and fails precisely where it is needed: on pinhole failing pairs
at view distance $\geq 4$ it finds \corrFarPinhole. G-geo-a is limited by
availability: structure-from-motion succeeds on \geoAavail\ of \nRowsStage\
rows, failing outright on two scenes, and dense stereo leaves holes exactly
where the small failing objects are. G-geo-b removes the availability problem
and its depth is accurate---median relative error \geoBdepthErr\% against
multi-view stereo where both exist---but its point-map resolution cannot resolve
the failing objects, and it passes \gateGeoB\ of \nGeoPairs.

G-gs deserves a paragraph because it is the bridge that the entire
3D-representation family of Sec.~\ref{sec:related} relies on. The
reconstructions are good by their own measure: training-view PSNR is
\gsPSNRmin--\gsPSNRmax~dB and the rendered depth covers the target objects.
Yet the depth error against multi-view stereo ranges from \gsDepthBest\% on the
scene with the \emph{lowest} PSNR to \gsDepthWorst\% on the scene with among
the highest, and across the \gsCams\ cameras with a reference depth the
rendered surface is nearer than the stereo surface on every one. With
\nCamMin--\nCamMax\ wide-baseline views, semi-transparent Gaussians placed
anywhere along a ray reproduce the training images; PSNR rises and the surface
dissolves into floaters. A depth that is dense and confidently wrong is worse
than one that is missing: the four pairs G-geo-a passed all collapse under
G-gs, and \gateGeoC\ of \nGeoPairs\ pass.

\subsection{Candidates exist; the selection signal does not}
The union of all generators contains a correct candidate for \unionOracle\ of
the \nQualPairs\ qualifying failing pairs, so the information is present in
the candidate set. The same union produces a non-empty candidate on
\unionHalluc\ of the ground-truth-empty rows. Any selector must therefore
separate ``the object, elsewhere'' from ``something else, anywhere'' without
ground truth. We registered one, SeedVote: a candidate replaces the default
seed only if it is dissimilar to the default, agrees with at least one other
stream or with a geometry-based scale check, scores higher than the default on
a SAM~3 confidence signal, and does not overlap another object; on rows where
the default seed is empty, a stricter agreement count is required so that the
hallucinating union cannot fill empty rows. The vote replaced \nVoteReplaced\
of \nGtPresent\ seeds. On the failing set it was accurate: \voteFzeroRep\
replacements, \voteFzeroRec\ recovered, \voteFzeroBroke\ broken, paired
$\Delta\Jscore=\voteFzeroDJ$ (CI \voteFzeroCI). On the sound set it was not:
\voteNzeroRep\ replacements, \voteNzeroBroke\ of which broke a working pair,
$\Delta\Jscore=\voteNzeroDJ$ (CI \voteNzeroCI). The net benchmark effect is
\voteRule\ points. With an oracle choosing among the same candidates the
effect would be \voteOracle\ points and \voteOraclePairs\ pairs recovered.
Hallucination on the empty rows did not increase. The hemisphere rigs alone
show a net gain with no new failures ($\Delta\Jscore=\voteHemiDJ$, CI
\voteHemiCI); restricting the rule to them would be a post-hoc choice, and we
do not make it.

\subsection{The structural reason}
\label{sec:asymmetry}
The result generalises beyond our rule. There are \nFzero\ pairs that need
repair and \nNzero\ that must not be touched, a ratio of \asymmetry. A
selector that is right on most of the failing pairs and wrong on a small
fraction of the sound ones still loses, because the small fraction is taken of
a population fifteen times larger; our own vote had a replacement precision of
83\% on F0 and 31\% on N0 and lost. Three independent attempts in this
project---a post-hoc seed repair, exemplar detection, and the vote---hit the
same wall. The quantity that determines whether seed repair on rigs can help is
not recall on the failing pairs, which every generator family can reach, but
precision on the sound population, and we suggest that future work report it
on that population directly.

% ===========================================================================
\section{Limitations}
\label{sec:limits}
\begin{enumerate}
  \item We did not re-run the XMem baseline. Its numbers are transcribed from
        the benchmark paper, and the equivalence of the two pipelines'
        structure is argued from the paper's description, not measured.
  \item There is one baseline and one dataset. SAM~2, Cutie, DEVA, and the
        2026 uncalibrated promptable methods do not appear as columns because
        none has been evaluated on MUVOD; running them is a contribution in
        itself and outside the present scope.
  \item Cross-view memory (\gainXview) failed its first pre-registered test
        and passed a second, re-registered one at the benchmark's reference
        camera. The mechanism is the same; the reader should weigh the
        component accordingly.
  \item The crop gate was selected on the scored scenes. Sec.~\ref{sec:loso}
        bounds the optimism at \losoOmega\ points.
  \item MUVOD scores three cameras per scene, so the Stage-1 populations of
        Sec.~\ref{sec:negative} live on two non-reference cameras per scene.
  \item The \nFisheye\ fisheye scenes are scored but excluded from the
        geometry arms, and from MUVOD's own 3D benchmark; the geometry
        conclusions are for pinhole rigs.
  \item The reference camera of one scene (Dog) is determined up to two
        candidates; both are reported.
\end{enumerate}

% ===========================================================================
\section{Conclusion}
On the only public rig benchmark, a two-stage pipeline with a promptable
foundation model improves on the published baseline by \DeltaBasic\ points,
of which \gainOurs\ is ours. The remaining \totalGap\ points are mostly ordinary
mask quality on pairs that never fail, not resolution and not drift. The
visible failures---seeds placed on the wrong object---can be repaired by an
oracle but not, in five pre-registered attempts, by any ground-truth-free rule,
because \nFzero\ pairs to fix stand against \nNzero\ to protect. That
\asymmetry\ asymmetry is the number the next method for this task should be
designed against.

% ===========================================================================
\balance
\bibliographystyle{IEEEtran}
\bibliography{refs}

\end{document}
